\chapter{Системийн шаардлага}
\label{chapter3}
\section{Системийн дотоод үйл ажиллагаа}
%Шинээр бүтээх системийн ажиллагааны тухайн дэлгэрэнгүй бичнэ. Ерөнхий ажиллагаа ямар байх, ямар ямар хэрэглэгчийн эрх үүсгэх, тэдгээр нь системд юу хийж чадах, ерөнхий ажиллагааны дараалал ямар байх гэх мэтийг маш дэлгэрэнгүй бичнэ. Тайлбарыг баяжуулах зорилгоор нэмэлт зураг, хүснэгт ашиглаж болно (жишээлбэл, үйл ажиллагааны бүдүүвчийг бүрэн зураглалаар харуулах, мөн хүн, банк, сервер зэргийн зургийг байрлуулж, хоорондын урсгалыг холбож зурах гэх мэт.)

\section{Системийг ашиглах хэрэглэгчид}
\subsection{Системийн оролцогчид}
%Системд оролцогчдыг тодорхойлж үүргийн нэрийг онооно.
%Нэр тус бүрийн ард тайлбар бичнэ. Жишээлбэл:
%Хувь хүмүүс - Хэрэглүүрийг ашиглан өөрийн зөв хоололтод анхаарал тавихыг хүссэн хэн бүхэн.
%Сервер/хостинг үйлчилгээ үзүүлэгчид - Системийг байршуулан ажиллуулах сервер түрээслэх үйлчилгээ үзүүлэгч байгууллагууд.

\subsection{Системийн тоглогчид}
%Системийн тоглогчдыг тодорхойлж үүргийн нэрийг онооно.
%Нэр тус бүрийн ард тайлбар бичнэ. Жишээлбэл:
%Хэрэглэгч - Хэрэглүүрийг ашиглах хүн.

%-------------------------------------------------------------------------------
\section{Функцийн шаардлага}
%Системийн функцийн шаардлагыг бичнэ. Бичихдээ шаардлага бичих дүрмийг дагаарай.
%-------------------------------------------------------------------------------
\section{Функцийн бус шаардлага}
%Системийн функцийн бус шаардлагыг бичнэ. Бичихдээ Ian Sommerville -н ангиллыг ашиглах
%-------------------------------------------------------------------------------
\section{Юзкейс диаграм}
%Visual Paradigm дээр зурсан анхны юзкейс диаграмыг зураг болгон энд оруулна.
%-------------------------------------------------------------------------------
\section{Юзкейсийн тодорхойлолт}
%Юзкейсийн тодорхойлолтуудыг бүрэн оруулна.
%-------------------------------------------------------------------------------
\section{Дэвшилтэт юзкейс диаграм}
%Visual Paradigm дээр зурсан дэвшилтэт хэлбэрийн юзкейс диаграмыг зураг болгон энд оруулна.
%-------------------------------------------------------------------------------
\section{Дэвшилтэт юзкейсийн тодорхойлолт}
%Дэвшилтэт юзкейсийн тодорхойлолтуудыг бүрэн оруулна.
%-------------------------------------------------------------------------------
\section{Системийн прототип}
% draw.io, visual paradigm, figma гэх мэт хэрэгслийг ашиглан програмын дэлгэцийн бүх интерфэйсүүдийг бүрэн зурж оруулах;